% Problem record op_cb7cdf37ec9b9ec5. This TeX file is derived from
% database/problems_json/op_cb7cdf37ec9b9ec5.json by scripts/sync-tex.mjs; edit the JSON record
% and rerun the script rather than editing this file.
\section{Gaussian optimality of two-mode entanglement of formation}
\label{sec:op_cb7cdf37ec9b9ec5}

\paragraph{Problem.}
Does Gaussian entanglement of formation equal unrestricted entanglement of formation for every finite-energy two-mode Gaussian state?

Let $\rho_{AB}$ be any finite-energy two-mode Gaussian density operator, with one bosonic mode held by each party. Finite energy means finite total mean photon number. Define

\begin{equation}
\begin{gathered}
E_F(\rho_{AB}):=\inf_\mu\int S(\operatorname{Tr}_B|\psi\rangle\langle\psi|)\,\mu(d\psi),
\\
S(\omega):=-\operatorname{Tr}(\omega\log_2\omega).
\end{gathered}
\label{eq:ser10-statement-1}
\end{equation}

The infimum in Eq.~\eqref{eq:ser10-statement-1} is over probability measures on normalised pure vectors with barycentre $\rho_{AB}$. The Gaussian entanglement of formation $E_F^{\mathrm G}$ restricts the same infimum to pure Gaussian vectors.

The proposed equality is $E_F(\rho_{AB})=E_F^{\mathrm G}(\rho_{AB})$, without an exchange-symmetry assumption.

\subsection*{Status}

Solved

\subsection*{Source}

Wolf et al., Section IX, discuss equality of the Gaussian and unrestricted convex roofs \sourcecite{ref:ser10-1}{WGKWC04}. Adesso resolves the full two-mode instance in Eq. (19) \sourcecite{ref:ser10-3}{Ade26}.

\subsection*{Progress}

\begin{itemize}
  \item Gaussian decompositions give the covariance optimisation

\begin{equation}
\begin{gathered}
E_F(\rho_{AB})\leq E_F^{\mathrm G}(\rho_{AB})
=\inf_{\substack{\gamma\preceq V\\\gamma\text{ pure Gaussian covariance}}}
g\!\left(\frac{\sqrt{\det\gamma_A}-1}{2}\right),
\\ g(x):=(x+1)\log_2(x+1)-x\log_2x,
\end{gathered}
\label{eq:ser10-progress-1-1}
\end{equation}

In Eq.~\eqref{eq:ser10-progress-1-1}, $V$ is the covariance of $\rho_{AB}$ in vacuum-$I$ units and $\gamma_A$ is the local covariance block of the pure Gaussian component. Use $0\log_2 0:=0$. Proposition 1 of Wolf et al. establishes this restricted optimization. \sourcecite{ref:ser10-1}{WGKWC04}

  \item Giedke et al., Proposition 2, establish $E_F=E_F^{\mathrm G}$ for symmetric two-mode Gaussian states. \sourcecite{ref:ser10-2}{GWKWC03}

  \item Adesso proves $E_F=E_F^{\mathrm G}$ for every two-mode Gaussian state in Eq. (19). Theorem 1 applies an affine entanglement bound to every finite-energy pure state, including non-Gaussian states. Equations (18)--(19) average it over arbitrary decompositions. Supplement S2.4 treats infinite-dimensional limits and probability-measure decompositions. \sourcecite{ref:ser10-3}{Ade26}
\end{itemize}

\subsection*{References}

\begin{enumerate}[label={},leftmargin=4.8em,labelsep=0.5em]
  \footnotesize
  \item[\textup{[WGKWC04]}]\label{ref:ser10-1}
  M. M. Wolf, G. Giedke, O. Kr\"uger, R. F. Werner, and J. I. Cirac, "Gaussian Entanglement of Formation," \emph{Physical Review A} \textbf{69}, 052320 (2004). \href{https://doi.org/10.1103/PhysRevA.69.052320}{doi:10.1103/PhysRevA.69.052320}; \href{https://arxiv.org/abs/quant-ph/0306177}{arXiv:quant-ph/0306177}.

  \item[\textup{[GWKWC03]}]\label{ref:ser10-2}
  G. Giedke, M. M. Wolf, O. Kr\"uger, R. F. Werner, and J. I. Cirac, "Entanglement of Formation for Symmetric Gaussian States," \emph{Physical Review Letters} \textbf{91}, 107901 (2003). \href{https://doi.org/10.1103/PhysRevLett.91.107901}{doi:10.1103/PhysRevLett.91.107901}; \href{https://arxiv.org/abs/quant-ph/0304042}{arXiv:quant-ph/0304042}.

  \item[\textup{[Ade26]}]\label{ref:ser10-3}
  G. Adesso, "Optimality of Gaussian Entanglement of Formation," preprint (2026), version 2, 13 August 2026. \href{https://doi.org/10.48550/arXiv.2608.01909}{doi:10.48550/arXiv.2608.01909}; \href{https://arxiv.org/abs/2608.01909v2}{arXiv:2608.01909v2}.
\end{enumerate}

\subsection*{Comment}

The affirmative resolution is an arXiv preprint, version 2; no peer-reviewed publication was identified in the September 2026 audit. It determines the single-copy convex roof. Additivity across copies remains separate. The record \texttt{01M1HME78068MQY7E9KA81B7WX} asks the distinct question for non-bisymmetric states with at least three modes.

\subsection*{Field}

Quantum Resource Theory

\subsection*{Topic}

Gaussian quantum information; Entanglement measures

\subsection*{ID}

\texttt{op\_cb7cdf37ec9b9ec5}
